\documentclass[UTF8,fontset=none,11pt,a4paper]{ctexart}
\usepackage{ipara-handbook}
\usepackage{amsmath,booktabs,tabularx,hyperref}
\usetikzlibrary{arrows.meta,positioning,fit,calc}

\handbooksetup{X-B01 Privilege / Exception / System Call}{408 · Cross-Subject Bridge}{Trap · Saved State · Kernel Control · Return}

\tikzset{
  unode/.style={draw=iparaDeep,fill=iparaSoft,rounded corners=1.5mm,align=center,minimum width=2.5cm,minimum height=0.88cm,font=\small},
  knode/.style={draw=iparaGreen,fill=white,rounded corners=1.5mm,align=center,minimum width=2.5cm,minimum height=0.88cm,font=\small},
  hwnode/.style={draw=iparaDeep,double,double distance=1.2pt,fill=white,rounded corners=1.5mm,align=center,minimum width=2.5cm,minimum height=0.88cm,font=\small},
  warnnode/.style={draw=iparaAccent,fill=white,rounded corners=1.5mm,align=center,minimum width=2.5cm,minimum height=0.88cm,font=\small},
  mainflow/.style={-{Latex[length=2mm]},thick,draw=iparaDeep},
  retflow/.style={-{Latex[length=2mm]},thick,dashed,draw=iparaGreen}
}

\begin{document}
\handbooktitle{Privilege / Exception / System Call × OS Control}{用户执行怎样安全跨越特权边界，并在内核处理后恢复或终止}{408 · Cross-Subject X-B01}

\begin{corebox}{Mother Interface}
\[
\boxed{
\begin{aligned}
\text{User Execution (U-Mode)} &\xrightarrow{\text{Trap / Syscall / Fault}} \text{Hardware Mode Elevation \& Auto-Save} \\
&\xrightarrow{\text{Vector Table}} \text{Kernel Dispatcher \& Service} \\
&\xrightarrow{\text{Hardware Return (\code{IRET/SRET})}} \text{Resume User Mode}
\end{aligned}
}
\]
本 Bridge 专门负责连接“计算机组成原理（硬件特权级、中断异常硬件通路）”与“操作系统（系统调用分发、异常修复与控制权决策）”的交接契约；微体系结构流水线（CO-04）与进程调度算法（OS-01/02）分别回到各自专题。
\end{corebox}

\begin{methodbox}{先定义接口两侧的词}
\begin{itemize}
  \item \kw{CPU 特权级 (Privilege Level)}：硬件强制实施的操作权限约束（如 RISC-V 的 U-mode 与 S/M-mode，x86 的 Ring 3 与 Ring 0）；特权指令（如关中断、修改页表基址寄存器）只能在内核态执行。
  \item \kw{陷入指令 (Trap / Syscall Instruction)}：用户态程序主动产生受控异常以请求内核服务的指令（如 x86 \code{syscall/int 0x80}，ARM \code{SVC}，RISC-V \code{ECALL}）。
  \item \kw{硬件陷入序列 (Hardware Interrupt/Trap Sequence)}：CPU 按 ISA 规定自动完成的\textbf{最小可信入口动作}，通常包括记录返回/异常位置与原因、更新特权/中断状态，并把 PC 转移到受保护入口；具体是否自动换栈、压栈或切页表是\textbf{体系结构相关}，不能写成跨架构定律。
  \item \kw{Trap Frame / 中断现场}：为了让被打断执行能够透明恢复而保存的体系结构/软件状态记录，通常包含通用寄存器与必要返回状态。它可以位于内核栈，也可以由 OS 使用专门内存区保存；\textbf{“一定在内核栈”不是定义的一部分}。
\end{itemize}
\end{methodbox}

\begin{methodbox}{Owners 与职责边界}
\begin{itemize}
  \item \kw{左侧 Owner (CO-02/03 指令系统与 CPU 控制器)：}拥有特权状态位（PSW/CSR）、陷入硬件逻辑、异常中断向量表寻址与返回指令（\code{IRET/SRET}）。
  \item \kw{右侧 Owner (OS-00/01 操作系统基础与进程)：}拥有系统调用分发号表、参数合法性校验、内核功能实现、信号投递与进程上下文管理。
  \item \kw{本 Bridge Owns：}跨特权边界时的现场保存协议、参数传递约定与返回/重试/终止决策流。
\end{itemize}
\end{methodbox}

\section{核心机制：跨越特权边界的硬件与软件分工}

用户程序绝不能直接跳转进内核任意地址（否则将彻底破坏安全性）。硬件与 OS 达成以下严格分工：
\begin{enumerate}
  \item \textbf{硬件责任（不可被软件绕过）}：
  \begin{itemize}
    \item 强制检查特权级；识别非法指令并阻断；
    \item 按 ISA 规定记录\textbf{最小返回状态与 trap 原因}，并更新必要的特权/中断状态；
    \item 把控制流转移到由特权软件预先配置的可信入口，使用户不能借 trap 任意选择内核 PC；
    \item 在返回指令执行时，根据受保护状态恢复先前特权级与继续执行位置。
  \end{itemize}
  \item \textbf{内核/入口软件责任（特权软件决策）}：
  \begin{itemize}
    \item 保存硬件没有自动保存、但透明恢复仍需要的通用寄存器，形成软件定义的 Trap Frame；
    \item 若当前 ISA/OS 入口协议需要，\textbf{由软件建立内核栈、切换页表或其它内核运行环境}；
    \item 读取系统调用号（如 \code{a7}）或异常原因状态（如 RISC-V \code{scause}），并完成分流；
    \item 校验用户态传入的指针与缓冲区边界合法性；
    \item 执行具体内核服务（如文件读写、页表修复），决定继续、重试、阻塞、终止或调度；
    \item 将执行结果写入约定返回位置，恢复软件保存的现场，再执行体系结构返回指令（如 \code{SRET/IRET}）。
  \end{itemize}
\end{enumerate}

\begin{warnbox}{架构边界：RISC-V/xv6 是最好的反例}
RISC-V 的 U-mode \code{ecall} 进入 S-mode 后，\textbf{不会自动切换 \code{sp} 或 \code{satp}，也不会自动保存全部通用寄存器}。xv6 先在 \code{trampoline.S:uservec} 中把用户寄存器保存到 per-process trapframe，再从 trapframe 取 \code{kernel\_sp} 与 \code{kernel\_satp}，由软件切换栈和页表后进入 \code{usertrap()}。因此应把“模式切换”“栈切换”“页表切换”“进程切换”视为四条可独立变化的状态轴。
\end{warnbox}

\section{母模型：特权转换与系统调用分发流转拓扑}

\begin{center}
\begin{tikzpicture}[x=3.40cm,y=1.55cm]
  % User Mode
  \node[unode] (uapp) at (0, 0.7) {User App\\[-1pt]\scriptsize \code{read(fd, buf, n)}};
  \node[unode] (trap) at (0, -0.7) {Trap Instruction\\[-1pt]\scriptsize \code{syscall / ecall}};

  % Hardware Sequence
  \node[hwnode] (hw) at (1.3, 0) {Hardware Sequence\\[-1pt]\scriptsize Save PC/PSW $\to$ S-Mode};

  % Kernel Mode
  \node[knode] (disp) at (2.6, 0.7) {Syscall Dispatch\\[-1pt]\scriptsize \code{sys\_call\_table[ID]}};
  \node[knode] (ret) at (2.6, -0.7) {Hardware Return\\[-1pt]\scriptsize \code{iret / sret} $\to$ U-Mode};

  % Connections
  \draw[mainflow] (uapp) -- (trap);
  \draw[mainflow] (trap) -- (hw);
  \draw[mainflow] (hw) -- (disp);
  \draw[mainflow] (disp) -- (ret);
  \draw[retflow] (ret) -- (uapp);
\end{tikzpicture}
\end{center}

\section{三种特权入口全景对比}

\begin{center}
\small
\renewcommand{\arraystretch}{1.25}
\begin{tabularx}{\linewidth}{>{\bfseries}p{0.18\linewidth} X X X}
\toprule
入口类型 & 触发源与同步性 & 硬件保存的断点 PC & 内核处理与后续出口 \\
\midrule
系统调用 (Syscall) & 用户程序主动执行陷入指令；\textbf{同步事件} & 逻辑恢复点通常是 syscall/trap 指令之后；\textbf{硬件具体保存哪个 PC 由 ISA 决定}。例如 RISC-V \code{ecall} 的 \code{sepc} 指向 \code{ecall} 本身，xv6 软件再将保存的 EPC 加 4 & 执行服务并返回结果，继续向下执行 \\
内中断 / 异常 (Fault) & 当前指令执行引发（缺页/除零/越界）；\textbf{同步事件} & 保存位置与异常类别由 ISA 定义；可恢复 fault 的语义通常要求修复后给故障指令再次执行机会 & 修复后重试、转入信号/异常处理，或终止进程 \\
外中断 (Interrupt) & CPU 外部硬件设备或时钟信号；\textbf{异步事件} & ISA 保存能够恢复被打断执行流的 continuation point；不要跨架构死记“永远等于 PC+某常数” & 处理外部事件，可能触发调度或恢复原流 \\
\bottomrule
\end{tabularx}
\end{center}

\section{最小完整执行轨迹：\texttt{read(fd, buf, n)} 跨界流转}

\begin{enumerate}
  \item \textbf{用户态发起}：用户 stub 按 ABI 把系统调用号写入约定寄存器（xv6/RISC-V 为 \code{a7}），参数位于 \code{a0, a1, a2,\ldots}，执行 \code{ecall}；
  \item \textbf{硬件最小跨界}：RISC-V 记录 trap 所需 CSR 状态（包括 \code{sepc/scause/sstatus}）、进入 S-Mode，并令 PC 跳到 \code{stvec}。\textbf{此时 xv6 仍使用 user page table 与用户寄存器现场}；
  \item \textbf{入口软件补全环境}：\code{trampoline.S:uservec} 把用户通用寄存器保存到 per-process trapframe，从中取出 \code{kernel\_sp/kernel\_satp}，由软件切到内核栈与内核页表，再进入 \code{usertrap()}；
  \item \textbf{原因判断与分发}：\code{usertrap()} 识别 \code{scause==8}。由于 RISC-V 此时 \code{sepc} 指向 \code{ecall} 本身，xv6 将保存的 EPC 加 4；随后 \code{syscall()} 从 trapframe 的 \code{a7} 读取服务号并通过 \code{syscalls[]} 分发到 \code{sys\_read()}；
  \item \textbf{内核校验与执行}：从保存现场取得参数，必要时经 \code{copyin/copyout/copyinstr} 等路径检查/复制用户地址，执行具体内核服务；
  \item \textbf{返回值与返回}：结果写入 trapframe 的 \code{a0}。内核准备 \code{sstatus/sepc} 与下一次 trap 元数据，\code{userret} 软件切回用户页表并恢复寄存器，最后 \code{sret} 恢复 U-Mode 与用户 PC。
\end{enumerate}

\begin{warnbox}{Anti-Bridge：三个必须阻断的误区}
\begin{itemize}
  \item \(\text{特权级切换} \neq \text{进程上下文切换}\)：系统调用可以始终代表同一个线程执行；即使某个 OS 的入口软件临时切换了 kernel stack 或 page-table root，也不能据此推出换了进程。
  \item \(\text{特权级切换} \neq \text{栈/页表切换}\)：前者是 CPU 权限状态变化；后两者是否随入口发生取决于 ISA 与 OS 设计。xv6/RISC-V 正是“先硬件升到 S-mode，后软件切 stack/satp”。
  \item \(\text{异常 (Fault)} \neq \text{中断 (Interrupt)}\)：异常由当前执行同步触发，中断通常来自异步外部事件；精确保存哪个 PC 与怎样恢复必须按 ISA/异常类别判断，不能跨架构机械套固定偏移。
  \item \(\text{进入内核} \neq \text{必须立刻调度}\)：系统调用处理完毕若无阻塞且时间片未耗尽，可以直接原路返回原线程。
\end{itemize}
\end{warnbox}

\begin{corebox}{30 秒极速调用协议}
做题遇到“特权级跨越、异常处理与中断返回”时，严格按三步判定：
\[
\boxed{\text{入口触发（主动/异常/外设）}} \xrightarrow{\text{硬件保存断点与特权}} \boxed{\text{内核返回：重试、继续还是终止}}
\]
\end{corebox}

\end{document}
